\documentclass[tikz,border=8pt]{standalone}
\usepackage{ctex}
\usepackage{amssymb}
\usepackage{pgfplots}
\pgfplotsset{compat=1.18}
\usetikzlibrary{calc,positioning,arrows.meta}

\begin{document}
\begin{tikzpicture}

%% ===== 左图：凸函数（碗形）=====
\begin{scope}[xshift=0cm]

  % 坐标轴
  \draw[->, gray!60, thin] (-3.2, 0) -- (3.2, 0) node[right, font=\scriptsize] {$x$};
  \draw[->, gray!60, thin] (0, -0.3) -- (0, 5.0) node[above, font=\scriptsize] {$f(x)$};

  % 凸函数 f(x) = 0.4*x^2 + 0.5
  \draw[blue!70!black, thick, domain=-3.0:3.0, samples=60]
    plot (\x, {0.4*\x*\x + 0.5});
  \node[blue!70!black, font=\small] at (2.5, 3.2) {$f(x)$};

  % 取两点 x1=-2, x2=2, 画弦
  \coordinate (P1) at (-2.0, {0.4*4+0.5});   % (-2, 2.1)
  \coordinate (P2) at (2.0, {0.4*4+0.5});    % (2,  2.1)

  \fill[black] (P1) circle (2pt) node[above left, font=\scriptsize] {$(x_1,f(x_1))$};
  \fill[black] (P2) circle (2pt) node[above right, font=\scriptsize] {$(x_2,f(x_2))$};

  % 弦（线段）
  \draw[red!80!black, thick] (P1) -- (P2);

  % 弦在曲线上方的阴影区域
  \fill[red!10, opacity=0.6]
    (-2.0, {0.4*4+0.5}) -- (2.0, {0.4*4+0.5})
    -- plot[domain=2.0:-2.0, samples=40] (\x, {0.4*\x*\x + 0.5}) -- cycle;

  % 弦在上方的标注
  \node[red!70!black, font=\scriptsize] at (0, 2.6) {弦在曲线上方};

  % 中点处的比较
  \coordinate (Mid) at (0, 2.1);     % 弦中点
  \coordinate (Fmid) at (0, 0.5);    % 函数值 f(0)=0.5
  \draw[<->, gray, thin] (0, 0.55) -- (0, 2.05) node[midway, right, font=\scriptsize] {Jensen 差};

  % 全局最小值
  \fill[red!60!black] (0, 0.5) circle (2pt)
    node[below right, font=\scriptsize] {全局极小 $x^*$};

  % 切平面下界示意（虚线切线）
  \draw[green!60!black, dashed, domain=-2.5:2.5]
    plot (\x, {0.5 + 0});   % 切线在最低点处是水平线
  \node[green!60!black, font=\scriptsize] at (2.5, 0.8) {切线（下界）};

  % 注释
  \node[draw=gray!50, fill=white, thin, rounded corners=2pt,
        font=\scriptsize, inner sep=3pt] at (0, -0.9)
    {$f(\theta x_1+(1-\theta)x_2)\leq\theta f(x_1)+(1-\theta)f(x_2)$};

  % 局部极小 = 全局极小
  \node[font=\scriptsize, blue!70!black] at (0, 4.5) {局部极小 $=$ 全局极小};

  % 标题
  \node[font=\small\bfseries] at (0, -1.8) {凸函数};

\end{scope}

%% ===== 右图：非凸函数（多峰）=====
\begin{scope}[xshift=8.0cm]

  % 坐标轴
  \draw[->, gray!60, thin] (-3.2, 0) -- (3.2, 0) node[right, font=\scriptsize] {$x$};
  \draw[->, gray!60, thin] (0, -0.3) -- (0, 5.0) node[above, font=\scriptsize] {$g(x)$};

  % 非凸函数：g(x) = 0.2*x^4 - 1.2*x^2 + 2.0
  % 极小值在 x = ±sqrt(3) ~ ±1.732，极大值在 x=0
  \draw[blue!70!black, thick, domain=-2.8:2.8, samples=80]
    plot (\x, {0.2*\x*\x*\x*\x - 1.2*\x*\x + 2.2});
  \node[blue!70!black, font=\small] at (2.4, 3.5) {$g(x)$};

  % 全局极小值（近似 x=±1.732，g=-0.4）
  \fill[red!60!black] (1.73, {0.2*8.99-1.2*3+2.2}) circle (2.5pt)
    node[below right, font=\scriptsize] {全局极小};
  \fill[red!60!black] (-1.73, {0.2*8.99-1.2*3+2.2}) circle (2.5pt)
    node[below left, font=\scriptsize] {全局极小};

  % 局部极大（x=0，g=2.2）
  \fill[orange!80!black] (0, 2.2) circle (2.5pt)
    node[above right, font=\scriptsize] {局部极大（鞍点）};

  % 取两点（位于不同极小值附近），弦穿出函数
  \coordinate (Q1) at (-2.0, {0.2*16-1.2*4+2.2});  % (-2, 1.0)
  \coordinate (Q2) at (2.0, {0.2*16-1.2*4+2.2});   % (2,  1.0)

  \fill[black] (Q1) circle (2pt) node[above left, font=\scriptsize] {$x_1$};
  \fill[black] (Q2) circle (2pt) node[above right, font=\scriptsize] {$x_2$};

  % 弦（在函数图像下方的区间）
  \draw[red!80!black, thick] (Q1) -- (Q2);

  % 弦在函数图像下方处的示意
  \coordinate (BelowMid) at (0, 1.0);
  \fill[orange!60!black] (BelowMid) circle (2pt);
  \node[below, font=\scriptsize, orange!60!black] at (BelowMid) {弦在曲线下方};

  % 阴影（函数图像在弦上方的区域）
  \fill[blue!10, opacity=0.5]
    plot[domain=-2.0:2.0, samples=60] (\x, {0.2*\x*\x*\x*\x - 1.2*\x*\x + 2.2})
    -- (2.0, 1.0) -- (-2.0, 1.0) -- cycle;

  % 非凸验证说明
  \node[draw=gray!50, fill=white, thin, rounded corners=2pt,
        font=\scriptsize, inner sep=3pt] at (0, -0.9)
    {存在 $x_1,x_2$：$g(\frac{x_1+x_2}{2})>\frac{g(x_1)+g(x_2)}{2}$（非凸）};

  % 梯度下降可能的路径（两条虚线箭头）
  \draw[->, thick, green!60!black, dashed, >=Stealth]
    (2.5, {0.2*39.0625-1.2*6.25+2.2}) -- (1.73, 0.4);
  \draw[->, thick, green!60!black, dashed, >=Stealth]
    (-2.5, {0.2*39.0625-1.2*6.25+2.2}) -- (-1.73, 0.4);
  \node[green!50!black, font=\scriptsize] at (-2.0, 3.8) {不同起点};
  \node[green!50!black, font=\scriptsize] at (-2.0, 3.3) {$\Rightarrow$不同极小};

  % 标题
  \node[font=\small\bfseries] at (0, -1.8) {非凸函数（多峰，有鞍点）};

\end{scope}

%% 整体标题
\node[font=\large\bfseries] at (7.8, 5.5) {凸函数（左）vs 非凸函数（右）};

\end{tikzpicture}
\end{document}
